% =============================================================================
%  LIMINAIRES, dédicace, remerciements, sigles, résumé, abstract, listes
%  (La première de couverture et la page de garde ne sont pas incluses.)
% =============================================================================

% --------------------------------------------------------------------- DÉDICACE
\entete{Dédicace}
\chapter*{DÉDICACE}
\addcontentsline{toc}{chapter}{DÉDICACE}

\vspace*{\fill}
\begin{center}
\begin{tcolorbox}[enhanced, nobeforeafter,
    colback=keycesoft, colframe=keyceaccent,
    boxrule=1.2pt, arc=10pt,
    width=0.78\textwidth,
    left=16pt, right=16pt, top=22pt, bottom=22pt,
    % Un second filet, plus fin et plus clair, à l'intérieur du premier :
    % le cadre a ainsi la sobriété d'un faire-part.
    borderline={0.4pt}{5pt}{keyceaccent!55}]
  \setstretch{1.15}%
  \centering
  {\large\itshape\color{keycegrey} Je dédie ce rapport}\\[12pt]
  {\large\itshape\color{keycegrey} à}\\[12pt]
  {\fontsize{19}{23}\selectfont\bfseries\color{keyceblue} la Famille KAMGANG}\\[16pt]
  {\color{keyceaccent}\rule{0.42\textwidth}{0.8pt}}\\[16pt]
  {\itshape\color{keycegrey} pour la confiance, la patience et les sacrifices
   consentis tout au long de ces années de formation.}
\end{tcolorbox}
\end{center}
\vspace*{\fill}

% ---------------------------------------------------------------- REMERCIEMENTS
\cleardoublepage
\entete{Remerciements}
\chapter*{REMERCIEMENTS}
\addcontentsline{toc}{chapter}{REMERCIEMENTS}

Un travail de fin d'études n'est jamais l'œuvre d'une seule personne. Nous
souhaitons exprimer notre reconnaissance envers toutes celles et tous ceux qui
ont joué un rôle dans l'élaboration de ce travail. C'est ainsi que nous
adressons un merci des plus sincères :

\begin{itemize}
  \item \textbf{aux membres du jury}, pour l'honneur qu'ils nous font en
        acceptant d'évaluer ce travail et pour le temps consacré à sa lecture ;

  \item \textbf{à M. GOUDJOU Blondon}, notre encadreur académique, pour ses
        précieux apports, son accompagnement constant et la rigueur d'esprit
        qu'il a su nous transmettre tout au long de la rédaction de ce travail ;

  \item \textbf{à M. Adrien GHOMSI}, notre encadreur professionnel au sein du
        \emph{Think Tank} Fou'ndjem, pour sa disponibilité, la pertinence de ses
        conseils et le suivi rigoureux qu'il a assuré durant le stage ; c'est à
        lui que nous devons la formulation initiale du problème de la visite
        immersive, dont ce travail est issu ;

  \item \textbf{à la Fondation Jean Félicien Gacha}, pour nous avoir accueilli
        en qualité de stagiaire, pour nous avoir ouvert ses espaces et ses
        archives, et pour nous avoir accordé la confiance ainsi que les
        ressources nécessaires à la réalisation de nos missions ;

  \item \textbf{à l'ensemble du corps enseignant de l'Institut Supérieur KEYCE
        Informatique \& Intelligence Artificielle}, campus de Bafoussam, pour
        la qualité des enseignements dispensés et pour leurs conseils ;

  \item \textbf{à tous les membres de notre famille}, qui n'ont cessé de nous
        épauler tant sur le plan moral que sur le plan financier ;

  \item \textbf{à nos camarades de promotion}, pour l'émulation, l'entraide et
        les relectures ;

  \item \textbf{à toutes les personnes} qui, n'étant pas citées plus haut, ont
        contribué d'une manière ou d'une autre à l'aboutissement de ce travail.
\end{itemize}

% -------------------------------------------------- LISTE DES SIGLES ET ABRÉVIATIONS
\cleardoublepage
\entete{Liste des sigles et abréviations}
\chapter*{LISTE DES SIGLES ET ABRÉVIATIONS}
\addcontentsline{toc}{chapter}{LISTE DES SIGLES ET ABRÉVIATIONS}

\begingroup
\setstretch{1.15}
\small
\begin{longtable}{@{}L{2.6cm}@{\hspace{0.6em}}L{11.6cm}@{}}
\textbf{ACM}      & \emph{AWS Certificate Manager}, service de gestion des certificats de sécurité \\
\textbf{API}      & \emph{Application Programming Interface}, interface de programmation applicative \\
\textbf{ARCore}   & Cadriciel de réalité augmentée de Google pour Android \\
\textbf{AWS}      & \emph{Amazon Web Services}, fournisseur d'infrastructure en nuage \\
\textbf{CDN}      & \emph{Content Delivery Network}, réseau de diffusion de contenu \\
\textbf{CI/CD}    & \emph{Continuous Integration / Continuous Deployment}, intégration et déploiement continus \\
\textbf{CRUD}     & \emph{Create, Read, Update, Delete}, les quatre opérations de base sur une donnée \\
\textbf{DNS}      & \emph{Domain Name System}, système de noms de domaine \\
\textbf{ECR}      & \emph{Elastic Container Registry}, registre d'images de conteneurs d'AWS \\
\textbf{ECS}      & \emph{Elastic Container Service}, orchestrateur de conteneurs d'AWS \\
\textbf{EKS}      & \emph{Elastic Kubernetes Service}, service Kubernetes géré d'AWS \\
\textbf{ERP}      & \emph{Enterprise Resource Planning}, progiciel de gestion intégré \\
\textbf{FCFA}     & Franc de la Coopération Financière en Afrique Centrale \\
\textbf{FinOps}   & \emph{Financial Operations}, discipline de pilotage des coûts du nuage \\
\textbf{GLB}      & \emph{GL Transmission Format Binary}, format de fichier 3D binaire \\
\textbf{HNSW}     & \emph{Hierarchical Navigable Small World}, structure d'index pour la recherche vectorielle \\
\textbf{HTTPS}    & \emph{HyperText Transfer Protocol Secure}, protocole web chiffré \\
\textbf{IA}       & Intelligence Artificielle \\
\textbf{IaC}      & \emph{Infrastructure as Code}, infrastructure décrite sous forme de code \\
\textbf{JSON}     & \emph{JavaScript Object Notation}, format d'échange de données \\
\textbf{JWT}      & \emph{JSON Web Token}, jeton d'authentification signé \\
\textbf{LiDAR}    & \emph{Light Detection And Ranging}, télémétrie par laser \\
\textbf{LLM}      & \emph{Large Language Model}, grand modèle de langue \\
\textbf{MIME}     & \emph{Multipurpose Internet Mail Extensions}, étiquette de type de fichier \\
\textbf{OIDC}     & \emph{OpenID Connect}, protocole d'authentification déléguée \\
\textbf{PWA}      & \emph{Progressive Web App}, application web installable \\
\textbf{QR}       & \emph{Quick Response}, code-barres à deux dimensions \\
\textbf{RA}       & Réalité Augmentée \\
\textbf{RAG}      & \emph{Retrieval-Augmented Generation}, génération augmentée par la recherche \\
\textbf{RLS}      & \emph{Row Level Security}, sécurité au niveau de la ligne \\
\textbf{S3}       & \emph{Simple Storage Service}, service de stockage d'objets d'AWS \\
\textbf{SaaS}     & \emph{Software as a Service}, logiciel en tant que service \\
\textbf{SGBD}     & Système de Gestion de Base de Données \\
\textbf{SPA}      & \emph{Single Page Application}, application web à page unique \\
\textbf{SQL}      & \emph{Structured Query Language}, langage de requête structuré \\
\textbf{TLS}      & \emph{Transport Layer Security}, protocole de chiffrement des échanges \\
\textbf{TTS}      & \emph{Text To Speech}, synthèse vocale \\
\textbf{URL}      & \emph{Uniform Resource Locator}, adresse web \\
\textbf{USDZ}     & \emph{Universal Scene Description Zipped}, format 3D d'Apple \\
\textbf{VR}       & \emph{Virtual Reality}, réalité virtuelle \\
\textbf{WebGL}    & \emph{Web Graphics Library}, interface graphique 3D du navigateur \\
\textbf{WebXR}    & Interface de réalité augmentée et virtuelle du navigateur \\
\end{longtable}
\endgroup

% ----------------------------------------------------------------------- RÉSUMÉ
\cleardoublepage
\entete{Résumé}
\chapter*{RÉSUMÉ}
\addcontentsline{toc}{chapter}{RÉSUMÉ}

\begingroup
\setstretch{1.15}

La valorisation numérique du patrimoine culturel constitue aujourd'hui un enjeu
majeur pour les institutions désireuses d'élargir leur audience au-delà de leurs
murs. La Fondation Jean Félicien Gacha, implantée à Bangoulap dans la région de
l'Ouest du Cameroun, a exprimé le besoin de permettre à ses visiteurs de
découvrir ses espaces sans contrainte de déplacement. C'est dans ce cadre que
s'inscrit notre stage, effectué du 20 juin au 20 août 2026 au sein du
\emph{Think Tank} Fou'ndjem de la Fondation. Le constat de départ est simple : la
découverte des espaces et des collections demeure conditionnée à une venue sur
place, ce qui écarte les publics éloignés, la diaspora et les personnes à
mobilité réduite. Lorsqu'une présence en ligne existe, elle se réduit le plus
souvent à une vitrine figée : des photographies plates, des notices courtes,
aucune sensation d'espace, aucun lien entre les objets et aucune possibilité de
poser une question.

L'objectif général de l'étude était de concevoir et de réaliser une plateforme
numérique unifiée capable de restituer à distance l'expérience de la visite.
L'étude relève du niveau intégrateur : il s'agit d'une recherche interactive,
conduite selon une approche mixte associant l'analyse documentaire,
l'observation directe des espaces, des entretiens semi-directifs avec le
personnel de la Fondation et une enquête par questionnaire auprès d'un
échantillon de visiteurs. La construction du logiciel a été menée de manière
itérative, selon une démarche DevOps où chaque modification du code est vérifiée
puis mise en production automatiquement.

La plateforme obtenue, nommée MUSÉA, articule quatre ensembles. Un progiciel de
gestion intégré, couplé à un site public, centralise la gestion des musées, des
salles, des œuvres, de la billetterie, de la boutique et des visiteurs. Un
parcours immersif à 360 degrés, complété par la modélisation en trois dimensions
et la réalité augmentée, permet de se déplacer dans les espaces ouverts au
public, de saisir un objet et de le poser à sa taille réelle chez soi. Un
assistant conversationnel fondé sur une architecture de génération augmentée par
la recherche répond aux questions des visiteurs à partir des seuls documents
validés par la Fondation. Enfin, un dispositif de rapprochement documentaire
retrouve, dans les collections ouvertes du monde, les œuvres apparentées à
celles de la Fondation. L'ensemble est déployé sur une infrastructure en nuage
dont le coût de fonctionnement est mesuré selon les principes du FinOps. Les
hypothèses de départ ont été vérifiées : la restitution spatiale, l'ancrage
documentaire de l'assistant et l'ouverture de l'architecture aux institutions
partenaires ont été observés et mesurés.

\vspace{0.5cm}
\noindent\textbf{Mots clés :} visite immersive ; réalité augmentée ;
génération augmentée par la recherche ; DevOps ; patrimoine numérique.
\endgroup

% --------------------------------------------------------------------- ABSTRACT
\cleardoublepage
\entete{Abstract}
\chapter*{ABSTRACT}
\addcontentsline{toc}{chapter}{ABSTRACT}

\begingroup
\setstretch{1.15}
\selectlanguage{english}

The digital promotion of cultural heritage has become a major challenge for
institutions wishing to reach audiences beyond their physical premises. The Jean
Félicien Gacha Foundation, based in Bangoulap in the West Region of Cameroon,
expressed the need to let visitors discover its spaces without having to travel.
This is the context of our internship, carried out from 20 June to 20 August
2026 within the Foundation's Fou'ndjem Think Tank. The initial observation is
straightforward: discovering the spaces and the collections still requires being
physically present, which excludes distant audiences, the diaspora and people
with reduced mobility. Where an online presence does exist, it is usually a
static showcase: flat photographs, short captions, no sense of space, no
connection between objects and no way of asking a question.

The general objective of the study was to design and build a unified digital
platform able to restore the experience of a visit remotely. The study belongs
to the integrative level: it is interactive research, conducted through a mixed
approach combining documentary analysis, direct observation of the premises,
semi-structured interviews with the Foundation's staff and a questionnaire
survey of a sample of visitors. Software construction proceeded iteratively,
following a DevOps approach in which every change to the code is verified and
then released automatically.

The resulting platform, named MUSÉA, brings together four sets of features. An
enterprise resource planning system, coupled with a public website, centralises
the management of museums, rooms, artefacts, ticketing, the shop and visitors.
A 360-degree immersive tour, complemented by three-dimensional modelling and
augmented reality, makes it possible to move through the areas open to the
public, to select an object and to place it at its real size at home. A
conversational assistant based on a retrieval-augmented generation architecture
answers visitors' questions using only documents validated by the Foundation.
Finally, a documentary matching mechanism finds, in the world's open
collections, the artefacts related to those of the Foundation. The whole system
is deployed on a cloud infrastructure whose running cost is measured according
to FinOps principles. The initial hypotheses were verified: spatial restitution,
the documentary grounding of the assistant and the openness of the architecture
to partner institutions were all observed and measured.

\vspace{0.5cm}
\noindent\textbf{Keywords:} immersive tour; augmented reality;
retrieval-augmented generation; DevOps; digital heritage.

\selectlanguage{french}
\endgroup

% ------------------------------------------------------------ LISTE DES TABLEAUX
\cleardoublepage
\entete{Liste des tableaux}
\addcontentsline{toc}{chapter}{LISTE DES TABLEAUX}
% L'interligne est réduit APRÈS le titre : appliqué avant, il déplacerait le
% cadre de titre de quelques points et cette page ne serait plus alignée sur
% les autres pages liminaires.
\listoftables
\vspace{-\baselineskip}

% ------------------------------------------------------------- LISTE DES FIGURES
\cleardoublepage
\entete{Liste des figures}
\addcontentsline{toc}{chapter}{LISTE DES FIGURES}
\listoffigures
\vspace{-\baselineskip}

% ---------------------------------------------------------------------- SOMMAIRE
\cleardoublepage
\entete{Sommaire}
\begingroup
\renewcommand{\contentsname}{SOMMAIRE}
\etocsetnexttocdepth{chapter}
\tableofcontents
\endgroup
